\chapter{Swollen glands}
\index{Swollen glands}

\medskip
\noindent\textit{Educational reference; always seek professional advice for individual care.}

\section*{Definition and Overview}
"Swollen glands" is the everyday term for lymphadenopathy, the enlargement of lymph nodes, the small, bean-shaped organs that filter fluid and fight infection throughout the body. Lymph nodes are part of the immune system, found in clusters in the neck, armpits, groin, behind the ears, under the jaw, and inside the chest and abdomen. When the immune system fights an infection, the nodes near the battle site swell as they produce infection-fighting cells, which is why a child with a sore throat often develops a tender lump under the jaw, or a person with an infected cut may feel a lump in the armpit. Most swollen glands are exactly this: a normal, temporary response to a minor infection, settling by themselves within a few weeks. However, swollen glands can occasionally signal something more serious, such as a chronic infection, an autoimmune disease, or cancer of the lymph system, so the challenge is to know when to reassure and when to investigate. This chapter explains what lymph nodes do, why they swell, how doctors decide whether a swollen gland is worrying, and what tests and treatments may be involved. The key message is reassuring but careful: most swollen glands are harmless and self-limiting, but a gland that is large, hard, fixed, persistent, or accompanied by other warning signs deserves proper medical assessment.

\section*{Epidemiology}
Swollen lymph nodes are extremely common, particularly in children, whose immune systems encounter new infections constantly. Studies of healthy children find that a substantial proportion have small, palpable lymph nodes in the neck at any given time, which come and go with colds and other minor illnesses. In adults, transient tender nodes are most often associated with throat, ear, dental, and skin infections, and they resolve as the infection clears. The pattern of swelling offers useful clues: a single localised cluster usually points to a problem in the area that the nodes drain, such as the neck nodes draining an infected throat or the armpit nodes draining an infected hand, while generalised swelling in several parts of the body suggests a whole-body process such as glandular fever, HIV, an autoimmune disease, or a blood cancer. Certain age groups have characteristic causes: young children mostly have viral and bacterial infections, adolescents and young adults are the typical age for glandular fever, and older adults are the group in whom a persistently enlarged node is more likely to require careful investigation for malignancy. Tuberculosis and HIV are major causes of lymphadenopathy in many regions, and although less common in many countries, they must be considered in at-risk people. Overall, the vast majority of swollen glands are benign, but the small proportion that are not makes persistence and pattern important signals.

\section*{Aetiology and Pathophysiology}
Lymph nodes act as filtering stations for lymph, the clear fluid that drains from the tissues. Each node contains specialised immune cells, particularly lymphocytes, which inspect everything that flows through for signs of infection or abnormal cells. A node swells when it is actively fighting an infection, when it is invaded by cells that do not belong there, such as cancer cells or certain organisms, or when a body-wide process causes inflammation. Infection typically produces tender, mobile, rubbery nodes that settle with the illness, whereas malignant involvement produces larger, harder, painless nodes that grow steadily.

\begin{itemize}
    \item \textbf{Viral infections:} the common cold, influenza, throat viruses, glandular fever, and childhood viruses such as measles all cause lymph node swelling, usually in the neck
    \item \textbf{Bacterial infections:} tonsillitis, strep throat, dental abscesses, ear infections, and infected wounds drain to nearby nodes, which become enlarged and tender
    \item \textbf{Glandular fever:} caused by the Epstein-Barr virus, it produces prominent, often bilateral neck nodes with a sore throat, fatigue, and sometimes fever
    \item \textbf{Tuberculosis:} causes chronic, matted, or painless lymph node enlargement, particularly in the neck, and needs specific testing and prolonged treatment
    \item \textbf{Other infections:} HIV, some sexually transmitted infections, toxoplasmosis, and various tropical and fungal infections can all cause lymphadenopathy
    \item \textbf{Autoimmune diseases:} conditions such as rheumatoid arthritis, lupus, and sarcoidosis cause generalised lymph node swelling through inflammation
    \item \textbf{Lymphoma and leukaemia:} cancers of the blood and lymph system classically cause painless, rubbery, or steadily enlarging nodes, often in several sites
    \item \textbf{Spread of other cancers:} cancer cells can travel to lymph nodes, such as breast cancer to the armpit, where they produce hard, fixed lumps
    \item \textbf{Medications:} a minority of drugs can cause node swelling as an allergic response
\end{itemize}

\section*{Clinical Features}
The features of swollen glands vary with the cause, and doctors pay attention to the details of the lump and its context.

\begin{itemize}
    \item \textbf{Lump in the neck, armpit, groin, or behind the ear:} the most obvious sign, ranging from a small pea-sized nodule to a large visible swelling
    \item \textbf{Tenderness:} common with infection and often the reason a person notices the node; painless swelling is more typical of chronic and malignant causes
    \item \textbf{Associated symptoms of infection:} sore throat, runny nose, ear pain, cough, or fever pointing to the source
    \item \textbf{Flu-like illness with fatigue:} typical of glandular fever, which is often mistaken for severe tonsillitis
    \item \textbf{Generalised swelling:} nodes in several areas at once, suggesting a body-wide process
    \item \textbf{Redness or warmth over the node:} suggests bacterial infection of the node itself, which may form an abscess
    \item \textbf{Night sweats, unexplained weight loss, or persistent fever:} red-flag symptoms that warrant prompt investigation
    \item \textbf{Steady growth or a node that stays enlarged for weeks:} a feature that always requires assessment, especially in adults
\end{itemize}

\section*{Differential Diagnosis}
Because many conditions present with swollen glands, the assessment is a detective process guided by location, size, consistency, duration, and accompanying symptoms.

\begin{itemize}
    \item \textbf{Common upper respiratory infections:} the commonest cause, easily recognised by the associated cold or throat symptoms and the self-limiting course
    \item \textbf{Glandular fever:} mimics strep throat but lasts longer, and is confirmed by a blood test for the Epstein-Barr virus
    \item \textbf{Salivary gland problems:} swelling of the parotid or submandibular glands, from infection, stones, or mumps, can be mistaken for lymph nodes
    \item \textbf{Thyroid enlargement or nodules:} a lump in the lower neck may be the thyroid rather than a lymph node, and needs ultrasound to clarify
    \item \textbf{Lymphoma and leukaemia:} persistent, painless, or enlarging nodes, especially with weight loss and night sweats, require biopsy or blood investigations
    \item \textbf{Tuberculosis:} chronic neck node enlargement in at-risk individuals, which may break down and drain
    \item \textbf{Connective tissue diseases:} generalised lymphadenopathy with joint pains, rash, or fever may be sarcoidosis, lupus, or rheumatoid arthritis
\end{itemize}

\section*{When to Seek Medical Care}
Most swollen glands need no medical care, but it is sensible to see a GP when a gland has not settled after two to four weeks, when it is growing, when it is hard or fixed to the skin, or when it is accompanied by persistent fever, weight loss, or night sweats. Children with swollen glands and high fever, difficulty breathing or swallowing, or a painful red swelling that suggests an abscess should be seen promptly. Anyone with a known cancer, or a family history of lymphoma or leukaemia, should have any new lymph node swelling assessed. There is no harm in a GP visit for reassurance, and the earlier a worrying cause is found, the better the outcome.

\textbf{Contact your local emergency services for:}
\begin{itemize}
    \item Difficulty breathing or swallowing with swollen neck glands
    \item Sudden difficulty breathing, wheezing, or swelling of the lips, tongue, or throat, which may indicate a severe allergic reaction
    \item High fever with a rapidly spreading tender red area in the neck
    \item Collapse or altered level of consciousness
\end{itemize}

\section*{Diagnosis and Investigations}
Assessment starts with a careful history and examination, and most cases need no tests at all.

\begin{itemize}
    \item \textbf{History:} the doctor asks how long the node has been present, whether it is growing, and about other symptoms, recent infections, travel, animal contact, and medicines
    \item \textbf{Examination:} the size, number, location, tenderness, texture, and mobility of the nodes are assessed, along with the areas they drain
    \item \textbf{Observation:} for typical infective nodes in a well child or adult, the plan is usually to wait two to four weeks and reassess
    \item \textbf{Blood tests:} a full blood count and inflammatory markers help distinguish infection from other causes, and specific tests for glandular fever, HIV, or autoimmune markers may be added
    \item \textbf{Ultrasound:} a quick, painless scan showing the structure of the node and whether it looks benign
    \item \textbf{Chest X-ray:} used if there are chest symptoms or suspicion of tuberculosis or lymphoma in the chest
    \item \textbf{Biopsy:} a sample of the node, taken with a needle or at operation, is the definitive test when malignancy, tuberculosis, or an unusual infection is suspected
    \item \textbf{Specialist referral:} persistent or worrying nodes are often referred to an ear, nose, and throat surgeon, haematologist, or infectious diseases physician
\end{itemize}

\section*{Management}
Treatment of swollen glands is treatment of the cause, since the node itself is only doing its job.

\begin{itemize}
    \item \textbf{Reassurance and observation:} the most common management for transient nodes from minor infection
    \item \textbf{Simple analgesia:} paracetamol or ibuprofen for pain and fever from an underlying infection
    \item \textbf{Antibiotics:} prescribed when a bacterial infection such as tonsillitis, strep throat, or a skin infection is diagnosed
    \item \textbf{Treatment of the source:} a dental abscess needs dental treatment, and an infected wound needs wound care, so the node can settle
    \item \textbf{Warm compresses and rest:} gentle home measures that ease tenderness while an infection resolves
    \item \textbf{Treatment of glandular fever:} no specific cure, but rest, fluids, and symptom relief, and avoidance of contact sports while the spleen may be enlarged
    \item \textbf{Treatment of specific infections:} tuberculosis needs a prolonged course of antituberculous drugs, and other infections have specific treatments
    \item \textbf{Treatment of lymphoma and leukaemia:} chemotherapy, immunotherapy, radiotherapy, or stem cell transplantation, coordinated by a specialist oncology team
    \item \textbf{Treatment of autoimmune disease:} disease-modifying medicines under specialist supervision
\end{itemize}

\section*{Complications}
\begin{itemize}
    \item \textbf{Abscess formation:} a node can become infected itself, forming a painful abscess that may need drainage
    \item \textbf{Delayed diagnosis of serious disease:} a persistently enlarged node that is not investigated can delay the diagnosis of lymphoma or tuberculosis
    \item \textbf{Splenic rupture in glandular fever:} a rare but dangerous complication, which is why contact sports are avoided for a month or more
    \item \textbf{Compression of surrounding structures:} very large nodes in the chest or neck can press on airways, blood vessels, or nerves
    \item \textbf{Chronic lymphoedema:} rarely, surgical removal of nodes for cancer treatment can impair lymph drainage and cause swelling of the limb
\end{itemize}

\section*{Prognosis and Outcomes}
The prognosis for swollen glands is almost always excellent, because the underlying cause is usually a short-lived infection. Nodes from viral or bacterial infections settle as the illness does, within a few days to weeks, and a residual small node may remain palpable for months before shrinking. Glandular fever takes longer, often several weeks, but is self-limiting and resolves fully in the vast majority of people. The important minority are cases caused by malignancy: the outlook depends on the specific cancer and its stage, but modern treatment of lymphoma is very effective, and early detection materially improves outcomes. For most people, the practical message is that a swollen gland is a sign of a working immune system, and the path to peace of mind is simple: allow a few weeks for an ordinary infection, and see a GP if the gland is large, hard, growing, or persisting, or warning symptoms appear. Reassurance from a careful assessment is itself a form of treatment, and most people never need more than that.

\section*{Prevention and Self-Care}
\begin{itemize}
    \item Practise good hygiene, including regular handwashing, to reduce the spread of the throat and skin infections that cause swollen glands
    \item Keep vaccinations, including influenza and measles, up to date
    \item Treat infected cuts and skin wounds promptly with cleaning and antiseptic, and see a doctor if they become increasingly red or tender
    \item Seek dental care early for toothaches, which can cause painful neck glands
    \item Rest, drink fluids, and use simple pain relief while an infective gland settles
    \item If a gland is growing, hard, painless, or present for more than a few weeks, see a GP rather than waiting
    \item Do not repeatedly squeeze or poke a swollen node, which can irritate it or introduce infection
\end{itemize}

\section*{Sources and Further Reading}
The following organisations provide reliable information about swollen glands and lymph node health.

\begin{itemize}
    \item healthdirect Australia. \textit{Swollen lymph nodes: what they mean.} https://www.healthdirect.gov.au/
    \item NHS. \textit{Swollen glands.} https://www.nhs.uk/
    \item Mayo Clinic. \textit{Swollen lymph nodes: symptoms and causes.} https://www.mayoclinic.org/
    \item Cancer Council Australia. \textit{Lymphoma: signs and symptoms.} https://www.cancer.org.au/
    \item Australian Department of Health and Aged Care. \textit{Immunisation and infection prevention.} https://www.health.gov.au/
    \item World Health Organization (WHO). \textit{Infectious disease prevention and control.} https://www.who.int/
\end{itemize}

\clearpage
